% =============================================
% Section 3: Problem Formulation
% =============================================
\section{Problem Formulation}
\label{sec:problem}

We formalize the constrained language model inference problem on
microcontrollers as a multi-objective optimization over model capacity,
deployment footprint, and language coherence.

\subsection{Hardware Constraint Model}

The \espss{} defines the following resource envelope:

\begin{table}[H]
\centering
\caption{Resource constraints of the \espss{} platform.}
\label{tab:constraints}
\small
\begin{tabular}{lrl}
\toprule
\textbf{Resource} & \textbf{Capacity} & \textbf{Role} \\
\midrule
Internal SRAM   & 512\,KB  & Stack, runtime variables \\
External \psram{} & 8\,MB & Model weights, KV-cache \\
Flash storage   & 16\,MB  & Firmware, model file, tokenizer \\
CPU frequency   & 240\,MHz & Inference computation \\
SIMD width      & 128-bit  & EE vector instructions \\
FPU             & None$^*$ & No hardware FP support \\
\bottomrule
\multicolumn{3}{l}{\footnotesize $^*$Software FP emulation available but extremely slow.}
\end{tabular}
\end{table}

Let $M_{\text{psram}} = 8$\,MB denote the \psram{} budget and
$M_{\text{flash}} = 16$\,MB the Flash budget (minus firmware overhead).
A deployed model must satisfy:

\begin{equation}
\underbrace{|\theta|_q}_{\text{quantized weights}}
+ \underbrace{S_{\text{kv}}}_{\text{KV-cache}}
+ \underbrace{S_{\text{act}}}_{\text{activations}}
\leq M_{\text{psram}}
\label{eq:memory_constraint}
\end{equation}

\noindent where $|\theta|_q$ is the size of the quantized model parameters,
$S_{\text{kv}} = 2 \cdot n_{\text{layers}} \cdot n_{\text{kv\_heads}} \cdot
d_{\text{head}} \cdot L \cdot b$ is the KV-cache size for sequence length $L$
and precision $b$ bytes per element, and $S_{\text{act}}$ accounts for
intermediate activations during inference.

\subsection{Language Coherence Metric}

We define a composite coherence score $\mathcal{C}$ that captures multiple
dimensions of language quality:

\begin{equation}
\mathcal{C}(m) = w_1 \cdot \hat{P}(m) + w_2 \cdot \text{MAUVE}(m) + w_3 \cdot \text{Judge}(m)
\label{eq:coherence}
\end{equation}

\noindent where:
\begin{itemize}[leftmargin=*,topsep=2pt,itemsep=1pt]
  \item $\hat{P}(m) = 1 - \frac{\log \text{PPL}(m)}{\log \text{PPL}_{\max}}$
        is the normalized inverse log-perplexity, scaled so that lower
        perplexity yields higher scores;
  \item $\text{MAUVE}(m) \in [0,1]$ measures distributional similarity to
        human-generated text~\cite{pillutla2021mauve};
  \item $\text{Judge}(m) \in [0,1]$ is the normalized LLM-as-judge score,
        following the TinyStories evaluation protocol adapted with
        GPT-4o~\cite{eldan2023tinystories};
  \item $w_1, w_2, w_3$ are weighting coefficients summing to 1.
\end{itemize}

\subsection{Optimization Objective}

The \tinybit{} design problem can be stated as:

\begin{equation}
\begin{aligned}
\max_{\theta, q} \quad & \mathcal{C}(\theta, q) \\
\text{s.t.} \quad & |\theta|_q + S_{\text{kv}} + S_{\text{act}} \leq M_{\text{psram}} \\
& |\theta|_q + S_{\text{tok}} + S_{\text{fw}} \leq M_{\text{flash}} \\
& \text{Latency}(\theta, q) \leq T_{\max}
\end{aligned}
\label{eq:optimization}
\end{equation}

\noindent where $q \in \{\text{FP32}, \text{INT8}, \text{INT4}, \text{Ternary}\}$
is the quantization scheme, $S_{\text{tok}}$ is the tokenizer size,
$S_{\text{fw}}$ is the firmware size, and $T_{\max}$ is the maximum acceptable
latency per token (we target $\leq 1$\,s, i.e., $\geq 1$ token/s).

This formulation makes explicit that the problem is not simply ``make the model
smaller,'' but rather finding the optimal point on the Pareto frontier of
model capacity, quantization aggressiveness, and coherence quality, all within
the hard constraints imposed by the microcontroller hardware.

\subsection{Ternary Weight Representation}

For the Xtensa kernel design (Section~\ref{sec:kernel}), we adopt the ternary
weight encoding of BitNet b1.58~\cite{ma2024bitnet}:

\begin{equation}
W_{ij} \in \{-1, 0, +1\}, \quad \text{encoded as 2 bits: } \begin{cases}
00 \rightarrow 0 \\
01 \rightarrow +1 \\
10 \rightarrow -1
\end{cases}
\label{eq:ternary}
\end{equation}

This encoding enables the \gemv{} operation $\mathbf{y} = W\mathbf{x}$ to be
decomposed into purely bitwise operations:

\begin{equation}
y_i = \sum_{j} W_{ij} \cdot x_j
    = \sum_{j \in \mathcal{P}_i} x_j - \sum_{j \in \mathcal{N}_i} x_j
\label{eq:bitwise_gemv}
\end{equation}

\noindent where $\mathcal{P}_i = \{j : W_{ij} = +1\}$ and
$\mathcal{N}_i = \{j : W_{ij} = -1\}$, eliminating all multiplication
operations.
